%% Copernicus Publications LaTeX template
%% Requires copernicus.cls (download from https://publications.copernicus.org/for_authors/manuscript_preparation.html)
\documentclass[journal abbreviation, manuscript]{copernicus}

\begin{document}

\title{Automatic camera-path generation for three-dimensional visualisation
of atmospheric rivers}

% \Author{firstname}{lastname}
\Author[1]{Suraj}{Kumar}
\Author[1]{}{}

\affil[1]{Department of Applied Computer Science, Fulda University of
Applied Sciences, Fulda, Germany}

\correspondence{Suraj Kumar (email)}

\runningtitle{Automatic camera paths for atmospheric rivers}
\runningauthor{S. Kumar}

\received{}
\pubdiscuss{}
\revised{}
\accepted{}
\published{}

\firstpage{1}

\maketitle

%% ============================================================
\section{Methodology}
%% ============================================================

\subsection*{Scope}

This section describes the automatic camera-path generation component.
Two inputs are taken as given: atmospheric-river (AR) detection, provided
as binary feature masks by an existing detection model supplied by our
supervisor, and feature tracking, which links detected features across
forecast time steps and was developed by a project colleague. The work
described here consumes those outputs and produces camera sequences that
Met.3D can execute; camera-path design was carried out jointly with a
second project member.

\subsection{Problem formulation}

Given a time-varying atmospheric field and a set of detected, tracked
atmospheric-river features, the task is to synthesise a camera trajectory
that presents those features intelligibly as a three-dimensional
animation. A camera trajectory is expressed as an ordered sequence of
keyframes, each specifying camera position in geographic coordinates
(longitude, latitude), a zoom/altitude parameter $z$, orientation angles
(pitch, yaw, roll), an interpolation flag, and a flag controlling whether
the underlying forecast time step advances at that key. Met.3D
interpolates between successive keys under a global tension parameter.

The design problem is therefore not merely geometric. A trajectory must
simultaneously (i) keep the feature of interest within frame, (ii)
preserve geographic context so that a viewer can locate the phenomenon,
(iii) remain smooth enough to be watchable, and (iv) coordinate camera
motion with the temporal evolution of the data.

\subsection{Data and pre-processing}

Detection output is provided as a NetCDF dataset on a
$1152 \times 768$ regular latitude--longitude grid with 61 forecast time
steps at three-hourly resolution, spanning a 180\,h forecast window.
Alongside the binary AR mask, the dataset provides total column water
vapour (TQV), 850\,hPa wind components, and mean sea-level pressure.
Tracking output is provided separately as a per-feature record
containing, for each time step, the feature's centroid position, area,
and physical extent, together with summary descriptors including
lifetime, total distance travelled, shape volatility, and the number of
split/merge events.

A single global time step typically contains many spatially disjoint AR
features. Individual features are therefore isolated by
connected-component labelling before any path is derived, with a
minimum-area threshold applied to reject noise.

\subsection{Centreline extraction}
\label{sec:centreline}

For an elongated feature, a natural camera trajectory follows its long
axis. Extracting that axis reliably required iterating over three
approaches.

\subsubsection{Morphological skeletonisation}

The initial approach applies morphological thinning to the binary mask,
reducing each feature to a single-pixel-wide centreline. This is
computationally inexpensive and is an established descriptor for
elongated atmospheric features. Two limitations emerged. First, skeletons
of irregular real-world shapes develop spurious branches; second, the
resulting centreline describes the geometric centre of the detected
shape, which need not coincide with the physical core of moisture
transport.

\subsubsection{Branch pruning by longest path}

To address branching, the skeleton is interpreted as an undirected graph
over 8-connected pixels, and its longest path (graph diameter) is
extracted by double breadth-first search. All pixels not on that path are
discarded. A useful side effect is that this procedure returns the
centreline already ordered end-to-end, which removes the need for a
separate coordinate sort.

This ordering property proved important. An earlier implementation
ordered centreline points by sorting on longitude. For features that
reverse direction in longitude, this interleaves points from different
parts of the feature and produces a path that repeatedly doubles back.
Walking the graph instead preserves true path order.

\subsubsection{Field-ridge centreline}

To address the geometric-versus-physical distinction, a third variant
extracts the centreline from the continuous moisture field rather than
the binary mask: sweeping across the feature's dominant axis, the cell of
maximum TQV is selected at each step. The resulting polyline follows the
ridge of moisture transport. This is a lightweight approximation of
topological ridge extraction via Morse--Smale complexes; a full
topological treatment was considered but not adopted, on the grounds that
the additional tooling cost was not justified for a trajectory that only
requires a smooth guide curve.

In all variants the extracted centreline is finally fitted with a
smoothing spline and resampled, yielding a continuous curve from which an
arbitrary number of waypoints can be drawn.

\subsection{Camera grammars}

Rather than a single trajectory type, several parameterised camera
behaviours (``grammars'') were implemented, each suited to a different
feature configuration:

\begin{itemize}
  \item \textit{Along-feature traversal.} The camera advances along the
        extracted centreline. Suited to strongly elongated features.
  \item \textit{Orbit.} The camera circles a feature on a ring of fixed
        angular radius, with orientation continuously directed inward.
        Suited to compact or highly volatile features, where following an
        erratic centroid would be visually unstable.
  \item \textit{Multi-feature tour.} Significant features are visited in
        sequence. Ordering by geographic longitude rather than by
        magnitude was adopted after the latter produced long
        intercontinental transits between consecutive targets.
  \item \textit{Regional cut.} Features are filtered to a geographic
        bounding box (for example the North Atlantic corridor) before a
        tour is constructed.
\end{itemize}

\subsection{Camera conventions and calibration}
\label{sec:conventions}

Met.3D's camera sequencer is only lightly documented, so the semantics of
the orientation parameters were established empirically by analysing
hand-authored sequences produced through the application's own interface.
Three findings materially changed the implementation.

\paragraph{Position is the camera's own location, not the target's.}
Hand-authored sequences contain latitude values outside the conventional
$\pm 90^{\circ}$ range. This indicates that the position fields describe
where the camera itself sits on its arc, which for oblique views lies
well away from the observed feature. An early implementation placed the
camera at the feature's own coordinates while raising pitch, which
produced views directed away from the data.

\paragraph{Angular offset should be bounded.}
Working sequences displace the camera from its target by a small fixed
angular distance, on the order of $5$--$15^{\circ}$. Scaling this offset
proportionally with pitch, as initially attempted, produces displacements
of comparable magnitude to the domain itself and moves the camera off the
region of interest.

\paragraph{Orientation is a bearing toward a target.}
Yaw is computed as the bearing from the camera position to the point
being observed, rather than from the direction of travel:
%
\begin{equation}
  \psi = -\arctan\!2\bigl(\lambda_{\mathrm{t}} - \lambda_{\mathrm{c}},\;
                          \phi_{\mathrm{t}} - \phi_{\mathrm{c}}\bigr),
  \label{eq:yaw}
\end{equation}
%
where $(\lambda_{\mathrm{c}}, \phi_{\mathrm{c}})$ is the camera position
and $(\lambda_{\mathrm{t}}, \phi_{\mathrm{t}})$ the target, in longitude
and latitude respectively. During traversal, the bearing is computed
toward the following waypoint rather than the current one, so that the
camera looks ahead along the feature. Yaw is additionally unwrapped
between consecutive keys so that rotation through the
$\pm 180^{\circ}$ discontinuity does not reverse direction.

\subsection{Temporal coordination}

Forecast time advances only on keys flagged for it. A key design
decision, adopted after comparison with hand-authored sequences, is to
separate camera motion from temporal motion: during segments in which the
camera moves, the forecast is held static; during segments in which the
forecast advances, the camera is held stationary. Sequences in which both
changed simultaneously were markedly harder to interpret.

Stationary segments are encoded as runs of identical keys with
interpolation disabled, with interpolation re-enabled on the final key of
the run to resume smooth motion. The number of such keys is budgeted
against the dataset length so that the total number of advances does not
exceed the available time steps.

\subsection{Composite sequence structure}

The resulting sequence template for a tracked feature is: a wide
establishing view; a descending approach toward the feature's current
position; a traversal or orbit of the feature; a withdrawal to an
intermediate altitude; and a stationary segment during which the forecast
advances and the feature evolves. For multi-cycle sequences this pattern
repeats, with each cycle re-approaching the feature at its updated
tracked position.

Continuity between segments is maintained explicitly. The camera state at
the end of each segment is retained and used as the initial state of the
next, so that no segment begins from a reset position. Only the opening
key originates from the global default view.

\subsection{Post-processing}

An optional smoothing stage operates on any generated sequence. Each
parameter is fitted with a cubic spline across the existing keys and
densely resampled, with angular parameters unwrapped before
interpolation. Optionally, orientation is re-derived from a look-ahead
point along the resampled path, which suppresses residual
high-frequency variation in heading. This stage is applied when the
underlying tracked centroid is itself volatile, trading a degree of
positional fidelity for visual stability.

%% ============================================================
\section{Results}
%% ============================================================

The pipeline was exercised on the 61-time-step forecast dataset described
above. Generated sequences load and execute in Met.3D, and produce
rendered animations of the detected atmospheric-river field. This section
reports the characteristics of generated sequences and the effect of the
principal design decisions.

\subsection{Quantitative characterisation of generated sequences}

Table~\ref{tab:sequences} summarises measured properties of hand-authored
reference sequences and of automatically generated output. Reported
quantities are the number of keys; the number of sharp direction
reversals (successive heading changes exceeding $150^{\circ}$), used as a
proxy for visual instability; the observed ranges of the zoom and pitch
parameters; and the number of keys advancing forecast time.

\begin{table}[t]
\caption{Measured properties of reference and generated camera
sequences.}
\label{tab:sequences}
\begin{tabular}{lccccc}
\tophline
Sequence & Keys & Rev. & $z$ range & Pitch range & Adv.\ keys \\
\middlehline
Reference: along-feature   & 26  & 0 & 73--250 & 0--44  & 25 \\
Reference: spline          & 84  & 1 & 16--250 & $-2$--66 & 64 \\
Reference: spin (single)   & 22  & 2 & 16--250 & 0--71  & 15 \\
Reference: spin+spline     & 64  & 1 & 25--250 & 0--52  & 0  \\
Reference: schema exemplar & 111 & 8 & 30--150 & 0      & 0  \\
Generated: multi-feature tour & 350 & 3 & 55--150 & 0   & 0  \\
\bottomhline
\end{tabular}
\belowtable{}
\end{table}

Two observations follow. First, hand-authored sequences employ a
substantially wider dynamic range in both zoom (16--250) and pitch (up to
$71^{\circ}$) than early automatically generated output, which was
confined to a conservative subset of the parameter space. Second,
reference sequences achieve low reversal counts with comparatively few
keys, indicating that Met.3D's own interpolation is responsible for much
of the observed smoothness, and that dense key sampling is not by itself
required.

\subsection{Effect of centreline ordering}

The ordering of centreline points was found to dominate trajectory
quality. A generated multi-feature sequence of 350 keys exhibited 80
sharp reversals when centreline points were ordered by longitude.
Replacing this with graph-order traversal of the pruned skeleton,
combined with ridge-based extraction, reduced the same measure to 3 --- a
reduction of approximately 96\,\%. A controlled test on a synthetic
U-shaped feature isolated the mechanism: longitude ordering produced 23
reversals where graph ordering produced none.

\subsection{Effect of camera-offset formulation}

Sequences generated with pitch-proportional angular offset placed the
camera at displacements exceeding $50^{\circ}$ from the observed feature,
and rendered views in which the feature left the frame partway through
the animation. Sequences generated with bounded offsets of
$5$--$15^{\circ}$, matching the convention observed in reference
material, retained the feature in frame throughout. Pitch magnitude alone
was not the determining factor: reference sequences use pitch up to
$71^{\circ}$ successfully when paired with appropriately bounded offsets.

\subsection{Grammar selection against feature descriptors}

Tracked features in the test dataset differ substantially in character,
supporting the case for multiple camera grammars rather than a single
trajectory type. The three highest-ranked tracked features each persisted
for the full 61-step window but differed markedly otherwise: total
distance travelled ranged from approximately 12\,400 to 27\,600\,km,
maximum extent from approximately 4000 to 11\,100\,km, and the number of
split/merge events from 11 to 59. The feature exhibiting 59 split/merge
events is poorly suited to centreline traversal, since its centroid is
unstable between consecutive time steps; an orbit about its mean position
produces a stable view of the same evolution.


%% ============================================================

\codeavailability{} %% TODO

\dataavailability{} %% TODO

\authorcontribution{} %% TODO

\competinginterests{The authors declare that they have no conflict of
interest.}

\begin{acknowledgements}
%% TODO
\end{acknowledgements}

\bibliographystyle{copernicus}
\bibliography{references.bib}

\end{document}
